\color{unidarkgreen}{\subsection{Логика деления стороны НН}}
\color{black}

\begin{enumerate}

\item
Логика деления стороны НН применяется в случае необходимости ступенчатого действия на отключение выключателей стороны НН (например, предварительное отключение выключателя ввода секции НН, затем отключение трансформатора от максимальной токовой защиты).

\item
Логика деления от МТЗ-1 НН1(НН) выводится накладкой \textbf{«ЛД НН1(НН) от МТЗ-1 НН1(НН)»}. Срабатывание логики деления формируется при возникновении пуска МТЗ-1 НН1(НН), далее, через выдержку времени \textbf{«Тср МТЗ-1 на В НН1(НН)»}, выдается сигнал срабатывания.

\item
Алгоритмы работы ЛД НН1(НН) от МТЗ-2, ЛД НН2 от МТЗ-1, ЛД НН2 от МТЗ-2, ЛД НН1(НН) от ТЗНП РЗН1 и ЛД НН2 от ТЗНП РЗН2 выполнены аналогично ЛД НН1(НН) от МТЗ-1.

\item
Алгоритм работы ЛД НН приведен на рисунке \ref{pic:LD_nn}. Уставки ЛД НН приведены в таблице \ref{app_set:ldnn}.

\medskip
\begin{figure}[H]
\centering
\includegraphics[width=1\textwidth,height=1\textheight,keepaspectratio]{\fbpath/06.01 ДЗТ/011.9902 ЛД НН/img/ld.pdf}
\captionof{figure}{Функциональная схема ЛД НН}\label{pic:LD_nn}
\end{figure}


\end{enumerate}